
%  Sección 7: Caracterización del Proceso
%  Subproceso 2: Selección y Registro de Cursos

\section{Caracterización del proceso}

El presente apartado describe los elementos fundamentales del Subproceso~2:
\textit{Selección y Registro de Cursos}, correspondiente al proceso de Matrícula
Regular del Estudiante en la Universidad Nacional de Ingeniería (UNI).

% ------------------------------------------------------------
\subsection{Actores}

Los actores son las personas, sistemas o áreas que participan directamente
en la ejecución del subproceso.

\begin{table}[H]
\centering
\caption{Actores del Subproceso~2}
\label{tab:actores_sp2}
\renewcommand{\arraystretch}{1.3}
\begin{tabular}{|p{3.8cm}|p{2.5cm}|p{7.5cm}|}
\hline
\textbf{Actor} & \textbf{Tipo} & \textbf{Rol} \\
\hline
Estudiante regular & Persona &
  Inicia el proceso, selecciona cursos y envía la propuesta de matrícula. \\
\hline
Intranet DIRCE & Software &
  Portal web (\texttt{dirce.uni.edu.pe}) donde el estudiante accede,
  visualiza la oferta y registra su selección. \\
\hline
SIGA-ORCE & Sistema &
  Motor \textit{backend} que valida restricciones (créditos, cruces,
  pre-requisitos) y gestiona las vacantes en tiempo real. \\
\hline
Comisión de Matrícula / Departamentos Académicos & Área institucional &
  Define la oferta habilitada, los turnos de acceso y los cupos
  por sección antes de iniciar el proceso. \\
\hline
\end{tabular}
\end{table}

% ------------------------------------------------------------
\subsection{Actividades}

Las actividades son las acciones que se realizan durante la ejecución
del subproceso, ordenadas secuencialmente.

\begin{table}[H]
\centering
\caption{Actividades del Subproceso~2}
\label{tab:actividades_sp2}
\renewcommand{\arraystretch}{1.3}
\begin{tabular}{|c|p{8.5cm}|p{3.5cm}|}
\hline
\textbf{N°} & \textbf{Actividad} & \textbf{Responsable} \\
\hline
1 & Acceder a la Intranet DIRCE en el turno habilitado. & Estudiante \\
\hline
2 & Visualizar la oferta de cursos y secciones habilitadas para el ciclo. & Sistema (SIGA-ORCE) \\
\hline
3 & Armar y seleccionar cursos, secciones y horarios (horario tentativo). & Estudiante \\
\hline
4 & Validar restricciones de forma interactiva: créditos, cruces de
    horario, pre-requisitos y disponibilidad de vacantes. & Sistema (SIGA-ORCE) \\
\hline
5 & Confirmar y enviar la propuesta de matrícula. & Estudiante \\
\hline
6 & Generar la Pre-Ficha de Matrícula y reservar vacantes temporalmente. & Sistema (SIGA-ORCE) \\
\hline
\end{tabular}
\end{table}

% ------------------------------------------------------------
\subsection{Recursos}

Los recursos son los medios necesarios para ejecutar el subproceso.

\begin{table}[H]
\centering
\caption{Recursos del Subproceso~2}
\label{tab:recursos_sp2}
\renewcommand{\arraystretch}{1.3}
\begin{tabular}{|p{4cm}|p{9.5cm}|}
\hline
\textbf{Recurso} & \textbf{Descripción} \\
\hline
Intranet DIRCE &
  Portal web institucional (\texttt{dirce.uni.edu.pe}); única vía de
  acceso al proceso de selección de cursos. \\
\hline
Servidores SIGA-ORCE &
  Infraestructura que procesa validaciones y gestiona el estado de
  vacantes en tiempo real. \\
\hline
Malla curricular digital &
  Catálogo de cursos con pre-requisitos, creditaje y secciones vigentes
  para el ciclo académico en curso. \\
\hline
Credenciales institucionales &
  Usuario y contraseña del estudiante para autenticarse en la intranet. \\
\hline
Reglamento RR-0570-2022 &
  Documento normativo que establece el límite de créditos a matricular
  según el promedio ponderado acumulado del estudiante. \\
\hline
\end{tabular}
\end{table}

% ------------------------------------------------------------
\subsection{Reglas}

Las reglas son las condiciones que deben cumplirse para que el
subproceso se ejecute correctamente.

\begin{table}[H]
\centering
\caption{Reglas de negocio del Subproceso~2}
\label{tab:reglas_sp2}
\renewcommand{\arraystretch}{1.3}
\begin{tabular}{|c|p{8cm}|p{3.3cm}|}
\hline
\textbf{ID} & \textbf{Regla} & \textbf{Fuente} \\
\hline
R1 & El estudiante solo puede acceder en el turno asignado por la DIRCE;
     fuera de él, el portal bloquea el registro de la selección. &
     Reglamento de Matrícula \\
\hline
R2 & El límite máximo de créditos a matricular depende del promedio
     ponderado acumulado del estudiante (véase Tabla~\ref{tab:creditos_sp2}). &
     RR-0570-2022 \\
\hline
R3 & El creditaje mínimo obligatorio por ciclo es de \textbf{12 créditos}. &
     Reglamento de Matrícula \\
\hline
R4 & Los cursos de arrastre (desaprobados o no llevados del ciclo previo)
     tienen prioridad en la selección. &
     Reglamento académico \\
\hline
R5 & Los pre-requisitos son obligatorios; el sistema bloquea
     automáticamente los cursos cuyos pre-requisitos no estén aprobados. &
     SIGA-ORCE \\
\hline
R6 & Los cruces de horario en laboratorios están \textbf{prohibidos};
     para cursos de teoría se toleran hasta \textbf{4 horas} de superposición. &
     Reglamento de horarios \\
\hline
R7 & Si no hay vacantes disponibles en los cursos requeridos, el
     estudiante debe solicitar rectificación de matrícula en el
     subproceso SP3A. &
     Proceso institucional \\
\hline
\end{tabular}
\end{table}

\begin{table}[H]
\centering
\caption{Límite de créditos según promedio ponderado acumulado (RR-0570-2022)}
\label{tab:creditos_sp2}
\renewcommand{\arraystretch}{1.3}
\begin{tabular}{|p{6cm}|c|}
\hline
\textbf{Promedio ponderado acumulado} & \textbf{Créditos máximos} \\
\hline
Menor a 11            & 18 \\
\hline
De 11,00 a 12,99      & 20 \\
\hline
De 13,00 a 14,99      & 22 \\
\hline
De 15,00 a 16,99      & 24 \\
\hline
De 17,00 a más        & 26 \\
\hline
\end{tabular}
\end{table}

% ------------------------------------------------------------
\subsection{Procedimientos del proceso}

Los procedimientos establecen la secuencia ordenada de pasos para
realizar las actividades del subproceso.

\begin{enumerate}
  \item \textbf{Habilitación del turno.} La DIRCE activa el acceso al
        portal para el estudiante según el cronograma de turnos publicado
        con anterioridad al inicio del periodo de matrícula.

  \item \textbf{Autenticación.} El estudiante ingresa con sus credenciales
        institucionales al portal \texttt{dirce.uni.edu.pe}.

  \item \textbf{Visualización de la oferta.} El sistema muestra los cursos
        disponibles para el ciclo, con sus secciones, horarios y número
        de vacantes disponibles en tiempo real.

  \item \textbf{Armado del horario tentativo.} El estudiante selecciona
        los cursos y secciones que desea cursar, priorizando los cursos
        de arrastre conforme al reglamento.

  \item \textbf{Validación automática.} SIGA-ORCE comprueba en tiempo real
        el cumplimiento de las reglas: (a)~límite de créditos según
        promedio, (b)~cruces de horario, (c)~pre-requisitos aprobados y
        (d)~disponibilidad de vacantes. Si alguna regla se incumple, el
        sistema alerta al estudiante y bloquea la acción hasta que se
        realice el ajuste correspondiente.

  \item \textbf{Envío de propuesta.} Una vez validada la selección, el
        estudiante confirma y hace clic en ``Registrar matrícula''.

  \item \textbf{Generación de Pre-Ficha.} El sistema genera la
        \textit{Pre-Ficha de Matrícula}, reserva temporalmente las
        vacantes de los cursos seleccionados y la registra en SIGA-ORCE.

  \item \textbf{Derivación al subproceso siguiente.} La Pre-Ficha de
        Matrícula generada constituye la entrada formal del
        Subproceso~3A: Pago y Formalización de la Matrícula. En caso de
        que algún curso no cuente con vacantes disponibles, el estudiante
        deberá gestionar una rectificación de matrícula a través del
        mismo subproceso SP3A.
\end{enumerate}
